\section{AI-Mediated Transformations and Mobile Lineage}
\label{sec:ai}

AI-assisted development is not one transformation. It ranges from completion of a few lines to generation of an application skeleton, selection of dependencies, modification across multiple files, test execution, debugging, and release preparation. The relevant provenance question is therefore not merely whether ``AI was used.'' It is which artifact-changing decisions were delegated, which context influenced them, and how the resulting outputs were bound to the final application.

\subsection{From text descriptions to mobile artifacts}

Text2App translates natural-language descriptions through an intermediate representation into Android application code~\cite{text2app}. This pipeline is useful for lineage analysis because it makes several transformations explicit: description to structured representation, representation to source, source to build inputs, and build inputs to APK. A generated application can share layout and behavior with other outputs while containing different source text. Conversely, two descriptions can produce similar source through templates or libraries.

AppForge is an Android development benchmark that evaluates generated projects and coding agents, including compilation-feedback refinement~\cite{appforge}. It supplies a functional evaluation setting, not an authenticated lineage system. In a tool-using workflow, prompts, retrieved material, tool responses, repository state, and accepted revisions can all affect the final output. A final source diff does not record that complete path.

These systems do not make traditional repackaging evidence obsolete. A generated app can itself be copied, repackaged, or combined with existing code. What changes is the space of legitimate relatedness. Multiple authorized outputs may descend from the same prompt, template, or model-mediated plan without one being copied from another. Direction cannot be inferred from pairwise similarity alone.

\subsection{Nondeterminism and reproducibility}

Conventional build reproducibility assumes that a declared input set and deterministic procedure should yield the same bytes. Generative systems may sample different outputs for the same prompt, call external tools, or depend on changing model and retrieval services. Exact reproduction can therefore be unavailable even when the process is authorized.

This does not eliminate reproducibility as an objective; it changes what must be reproduced. One target is \emph{record reproducibility}: an independent verifier can confirm the prompt or specification reference, model or service identity, parameters, tool sequence, retrieved materials, and accepted outputs. Another is \emph{decision replay}: a recorded output is retained and subsequent deterministic build steps reproduce the APK. A third is \emph{distributional replication}: repeated generations satisfy a defined property, which is an empirical model claim rather than artifact lineage.

For mobile supply-chain assurance, preserving the accepted generated source or patch is usually more important than reproducing the model's sampling event. The build attestation should treat that accepted output as a material with a digest. If policy requires regeneration, the record must acknowledge that regenerated code may be a sibling rather than an identical ancestor.

\subsection{Security of generated code}

Studies of code assistants show that generated suggestions can contain security weaknesses and that user interaction affects the result. Pearce et al. evaluated Copilot contributions across security-relevant scenarios~\cite{pearce}. Perry et al. found less secure code in their assisted group, whereas Sandoval et al. found a small security impact in a constrained C task~\cite{perry,sandoval}. Khoury et al. examined ChatGPT-generated code and improvement prompts~\cite{khoury}. The differing tasks, models, participants, and vulnerability oracles preclude a pooled claim that assistance invariably increases insecurity.

For lineage, the important finding is structural: generation provenance does not imply secure behavior. An authorized model-assisted change may introduce a vulnerability; an unauthorized derivative may remove one. Security review must analyze the actual accepted code and built artifact. The provenance record helps identify which model-mediated step introduced a change, but the vulnerability evidence comes from static analysis, testing, review, or formal reasoning.

Human approval also needs careful interpretation. Clicking ``accept'' establishes an interaction event, not necessarily informed review. A policy may require that a maintainer approve generated changes, but the evidential value depends on what was shown, which tests ran, whether warnings were visible, and whether the reviewer had authority for the affected component. Authorization and assurance remain separate.

\subsection{Dependencies and package hallucination}

Generative systems can recommend dependencies by name. Spracklen et al. show that code-generating models can produce nonexistent package names, creating a package-confusion risk when an attacker registers a hallucinated name~\cite{spracklen}. In an Android project, analogous risks can involve Maven coordinates, Gradle plugins, native packages, JavaScript dependencies in hybrid apps, or download URLs. The suggestion may enter source or build configuration even when no copied code appears in the model output.

This risk links AI assistance directly to supply-chain provenance. A robust record should capture the dependency request, resolver result, repository origin, version, integrity digest, and whether the dependency was already approved. A policy can reject unrecognized repositories or require lockfile review. Binary component analysis then checks whether the resolved material appears in the APK. Neither the prompt log nor the binary detector alone provides the full chain.

Package hallucination also creates a novel counter-history for library attribution. A newly registered malicious package may intentionally implement an interface or name suggested by a model. The package's code is not necessarily similar to an established library, so clone detection may not flag it. Repository provenance, registration chronology, maintainer identity, and dependency-policy evidence become central.

\subsection{Training-data and privacy lineage}

CodexLeaks demonstrates that code-generation models can emit sensitive information associated with training data~\cite{codexleaks}. The resulting lineage question differs from ordinary source copying. A generated fragment may be statistically or verbatim related to training material whose exact membership is unavailable. A final application can therefore contain sensitive or proprietary material without a normal repository-level copy event.

Artifact similarity can still help when a suspected source is known, but absence of a match is weak evidence because generation may transform names and structure. Provenance records can establish which model and prompt produced the accepted output, yet they may not reveal the relevant training example. Policy and legal conclusions require evidence beyond technical resemblance, including license, confidentiality, and permitted-use analysis.

Privacy also affects what provenance should retain. Storing complete prompts, retrieved files, and model outputs may itself expose secrets. A mobile-development provenance design should support selective disclosure: retain cryptographic commitments, access-controlled detailed records, and public metadata sufficient to bind the final artifact without publishing sensitive context. The record needs an evidence-retention policy, not indiscriminate logging.

\subsection{Multi-principal and tool-mediated authorship}

An agentic workflow can include a human requester, model provider, model instance, retrieval source, code repository, build service, package registry, test system, reviewer, and release signer. These are not equivalent principals. The model service may produce candidate text but lack release authority. A human may approve a patch but not control the signing key. A build service may emit an attestation but not own the source.

Supply-chain layouts are well suited to representing these roles when model-mediated steps are treated as explicit transformations. The layout can require that generated outputs pass tests, static checks, dependency policy, and human review before signing. It can also distinguish advisory output from accepted material. This is more precise than adding one boolean ``AI-generated'' field to the APK metadata.

The same role separation helps incident response. If a vulnerability is found, investigators can ask whether it originated in a generated patch, retrieved example, dependency recommendation, manual edit, or build-time transformation. The answer guides remediation and model or policy changes. Without step records, the final Git history may collapse several causes into one commit.

\begin{table}[t]
\caption{AI-mediated development risks and the evidence needed to reason about them.}
\label{tab:ai-risks}
\footnotesize
\begin{tabularx}{\textwidth}{L{2.35cm}Y Y L{1.6cm}}
\toprule
Risk or ambiguity & Why final-artifact similarity is insufficient & Complementary evidence & Main claim \\
\midrule
Sibling generations from one specification~\cite{text2app,appforge} & Similarity does not reveal whether one output copied another & Specification binding, generation records, accepted-output digests & \claimlevel{3}--\claimlevel{4} \\
Insecure generated code~\cite{pearce,perry,sandoval,khoury} & Origin says nothing about security of accepted behavior & Static/dynamic analysis, tests, review evidence, exact artifact binding & \claimlevel{6} \\
Hallucinated dependencies~\cite{spracklen} & Malicious package may be novel and not resemble a known library & Resolver provenance, repository identity, lockfile, component verification & \claimlevel{1}+\claimlevel{4} \\
Training-data or privacy leakage~\cite{codexleaks} & Relevant ancestor may be unknown or transformed & Model/prompt record, similarity investigation, access-controlled context & Conditional \claimlevel{3} \\
Agentic multi-file edits~\cite{appforge} & One final diff hides intermediate tool calls and rejected outputs & Step attestations, tool transcripts, policy checkpoints & \claimlevel{4}--\claimlevel{5} \\
Nondeterministic regeneration & Exact output may not be reproducible from a prompt alone & Preserve accepted output; separate record replay from model replication & \claimlevel{0}+\claimlevel{4} \\
\bottomrule
\end{tabularx}
\end{table}

\subsection{A minimum evidence record for generated transformations}

A generated transformation record should identify: the parent source revision; task or specification identifier; model or service identity at the granularity available; generation parameters; tool and retrieval accesses; candidate-output digests; the accepted patch or file digests; tests and analyses run; human or policy approvals; and the downstream build product. Sensitive context can be access controlled, but the public or shared record should still bind the accepted output to the final APK.

This record serves three purposes. It establishes direction from a parent revision to an accepted generated change. It separates generation from approval and release. It gives artifact and behavior analyses exact objects to examine. The record does not prove that the model never used unrecorded training material or that the accepted code is secure. It makes those residual questions explicit rather than hiding them behind a generic provenance claim.
